
\documentclass{.local/lib/classes/ndvr-vector-image-standalone}

\title[NDVR thesis-abstract]{%
    Темпоральная модель визуального потока видео%
}

\textwidth=40cm

\begin{document}

    %% My Types

    \newcommand{\umldtp}[1]{\textsf{#1}}
    \newcommand{\umlctp}[1]{\textsl{$\star$ #1}}
    \newcommand{\umlfld}[1]{\textsf{#1}}


    %% Tabular Styling

    \newcolumntype{F}[1]{>{\raggedleft\arraybackslash}m{#1}}
    \newcolumntype{T}[1]{>{\raggedright\arraybackslash}m{#1}}

    \setlength\dashlinedash{0.2pt}
    \setlength\dashlinegap{2.5pt}
    \setlength\arrayrulewidth{0.2pt}

    \renewcommand{\arraystretch}{1.2}
    \setlength{\tabcolsep}{0.1em}


    %% Tikz Uml Styling

    \colorlet{commonClass}{yellow!10}
    \colorlet{commonNote}{green!10}
    \colorlet{commonDraw}{lightgray}

    \colorlet{emphClass}{yellow!20}
    \colorlet{emphDraw}{black}

    \tikzumlset{
        %font=\normalsize,
        fill class=commonClass,
        fill note=commonNote,
        draw=commonDraw,
    }

    \tikzstyle{tikzuml inherit style}=[
        color=black,
        very thick,
        {-{Stealth[inset=0pt,scale=1.5,fill=white,angle'=45]}}
    ]

    \tikzstyle{tikzuml aggregation style}=[
        color=black,
        very thick,
        {{Diamond[scale=2,fill=white]}-}
    ]

    \tikzstyle{tikzuml composition style}=[
        color=black,
        very thick,
        {{Diamond[scale=2,fill=gray!90]}-}
    ]

    \tikzstyle{tikzuml implements style}=[
        color=black,
        very thick,
        {-{Stealth[inset=0pt,scale=1.5,fill=white,angle'=45]}}, 
        dashed
    ]



    \tikzstyle{videof} = [
        rectangle, rounded corners,
        ultra thick,
        minimum height=5em,
        minimum width=17em,
        draw=black!85,
        top color=black!5,
        bottom color=black!5
    ]

    \tikzstyle{scenef} = [
        rectangle, rounded corners,
        ultra thick,
        minimum height=5em,
        minimum width=8em,
        draw=black!85,
        top color=black!5,
        bottom color=black!5
    ]

    \tikzstyle{shotf} = [
        rectangle, rounded corners,
        ultra thick,
        minimum height=5em,
        minimum width=4em,
        draw=black!85,
        top color=black!5,
        bottom color=black!5
    ]

    \tikzstyle{frame1f} = [
        rectangle, rounded corners,
        ultra thick,
        minimum height=5em,
        minimum width=2em,
        draw=black!85,
        top color=black!5,
        bottom color=black!5
    ]

    \tikzstyle{frame2f} = [
        rectangle, rounded corners,
        ultra thick,
        minimum height=5em,
        minimum width=1em,
        draw=black!85,
        top color=black!5,
        bottom color=black!5
    ]

    \tikzstyle{layerf} = [
        rectangle,dashed,
        ultra thick,
        minimum width =1em,
        minimum height = 1em,
        inner sep=0.5em,
        draw=black!75,
        decorate,decoration={random steps,segment length=0.5em,amplitude=0.1em}
    ]

    \newcommand{\newSin}[1]
        {sin((#1)*0.5*pi)}

    \newcommand{\newXSin}[1]
        {\newSin{#1}*(#1)}

    \newcommand{\newXXSin}[1]
        {sin((#1)*2.0*pi)}

    \pgfooclass{stamp}{ % This is the class stamp
        \method stamp() { % The constructor
        }
        \method apply(#1) {
            %Draw the stamp:
            \draw[smooth,blue] plot[id=point1-shot1,domain=0:4,samples=10]
                function{0.3*\newXSin{x}};
            \draw[smooth,blue] plot[id=point1-shot2,domain=4:8,samples=10]
                function{-0.5*\newSin{x - 4}};
            \draw[smooth,blue] plot[id=point1-shot3,domain=8:16,samples=20]
                function{0.7*\newSin{x}};
            \draw[smooth,red] plot[id=point2-shot3,domain=8:12,samples=10]
                function{0.7*\newSin{x} - 0.4};
            \draw[smooth,red] plot[id=point2-shot4,domain=12:16,samples=10]
                function{0.5*\newXXSin{x} + 0.4};
        }
    }

    \pgfoonew \mystamp=new stamp()

    \pgfdeclarelayer{points}
    \pgfdeclarelayer{background}
    \pgfdeclarelayer{foreground}
    \pgfsetlayers{background,main,foreground,points,connections}



    %% Main Part

    \sf\singlespacing
    \begin{tikzpicture}[
        very thick,
        %node distance=8em
    ]
        \umlsimpleinterface{Видео}

        \umlclass[
            below=4em of Видео,
            draw=emphDraw,
            fill=emphClass,
        ]{Поток}{
            \begin{tabular}{crcl}
                $+$ & \umlfld{Номер} 
                    &:& \umldtp{Число};\\
                \hdashline
                $+$ & \umlfld{Время начала} 
                    &:& \umlctp{Точка времени или $\emptyset$};\\
                \hdashline
                $+$ & \umlfld{Время конца} 
                    &:& \umlctp{Точка времени или $\emptyset$};\\
                \hdashline
                $+$ & \umlfld{Опции} 
                    &:& \umlctp{Список опций};\\
                \hdashline
                $-$ & \umlfld{Сцены} 
                    &:& \umlctp{Последовательность сцен};\\
            \end{tabular}
        }

        \umlaggreg[mult1=1, mult2=1..$\ast$]{Видео}{Поток}

        \umlclass[
            below=4em of Поток,
            draw=emphDraw,
            fill=emphClass,
        ]{Сцена}{
            \begin{tabular}{crcl}
                $+$ & \umlfld{Номер} 
                    &:& \umldtp{Число};\\
                \hdashline
                $+$ & \umlfld{Время начала} 
                    &:& \umlctp{Точка времени};\\
                \hdashline
                $+$ & \umlfld{Время конца} 
                    &:& \umlctp{Точка времени или $\emptyset$};\\
                \hdashline
                $+$ & \umlfld{Опции} 
                    &:& \umlctp{Список опций};\\
                \hdashline
                $-$ & \umlfld{Съёмки} 
                    &:& \umlctp{Последовательность съёмок};\\
            \end{tabular}
        }

        \umlaggreg[mult1=1, mult2=1..$\ast$]{Поток}{Сцена}

        \umlclass[
            below=4em of Сцена,
            draw=emphDraw,
            fill=emphClass,
        ]{Съёмка}{
            \begin{tabular}{crcl}
                $+$ & \umlfld{Номер} 
                    &:& \umldtp{Число};\\
                \hdashline
                $+$ & \umlfld{Время начала} 
                    &:& \umlctp{Точка времени};\\
                \hdashline
                $+$ & \umlfld{Время конца} 
                    &:& \umlctp{Точка времени или $\emptyset$};\\
                \hdashline
                $+$ & \umlfld{Опции} 
                    &:& \umlctp{Список опций};\\
                \hdashline
                $-$ & \umlfld{Кадры} 
                    &:& \umlctp{Последовательность кадров};\\
            \end{tabular}
        }

        \umlaggreg[mult1=1, mult2=1..$\ast$]{Сцена}{Съёмка}


        \umlclass[
            below=4em of Съёмка,
            draw=emphDraw,
            fill=emphClass,
        ]{Кадр}{
            \begin{tabular}{crcl}
                $+$ & \umlfld{Номер} 
                    &:& \umldtp{Число};\\
                \hdashline
                $+$ & \umlfld{Время начала} 
                    &:& \umlctp{Точка времени};\\
                \hdashline
                $+$ & \umlfld{Время конца} 
                    &:& \umlctp{Точка времени};\\
                \hdashline
                $+$ & \umlfld{Опции} 
                    &:& \umlctp{Список опций};\\
                \hdashline
                $-$ & \umlfld{Данные} 
                    &:& \umldtp{Байтовый массив};\\
                \hdashline
            \end{tabular}
        }

        \umlaggreg[mult1=1, mult2=1..$\ast$]{Съёмка}{Кадр}


        \umlclass[
            right=3em of Поток,
        ]{Опция}{
            \begin{tabular}{cF{5em}cT{11em}}
                $+$ & \umlfld{Имя} 
                    &:& \umldtp{Строка};\\
                \hdashline
                $+$ & \umlfld{Тип} 
                    &:& \umldtp{\{{\tt\footnotesize  'Строка'}, {\tt\footnotesize  'Число'}, ...\}};\\
                \hdashline
                $+$ & \umlfld{Значение} 
                    &:& \umldtp{Байтовый массив};\\
            \end{tabular}
        }


        \umlaggreg[
            mult1=1, 
            mult2=0..$\ast$
        ]
            {Поток}{Опция}

        \umlaggreg[
            mult1=1, 
            mult2=0..$\ast$,
            geometry=-|
        ]
            {Сцена}{Опция}

        \umlaggreg[
            mult1=1, 
            mult2=0..$\ast$,
            geometry=-|
        ]
            {Съёмка}{Опция}

        \umlaggreg[
            mult1=1, 
            mult2=0..$\ast$,
            geometry=-|
        ]
            {Кадр}{Опция}

        \umlclass[
            left=3em of Поток,
        ]{Точка времени}{
            \begin{tabular}{cF{9em}cT{7em}}
                $+$ & \umlfld{Время отображения} 
                    &:& \umldtp{Число};\\
                \hdashline
                $+$ & \umlfld{Время кодирования} 
                    &:& \umldtp{Число};\\
                \hdashline
                $+$  & \umlfld{Масштаб времени}
                     &:& \umldtp{(Число, Число)};\\
            \end{tabular}
        }

        \umlaggreg[
            mult1=1, 
            mult2=0..2
        ]
            {Поток}{Точка времени}

        \umlaggreg[
            mult1=1, 
            mult2=1..2,
            geometry=-|
        ]
            {Сцена}{Точка времени}

        \umlaggreg[
            mult1=1, 
            mult2=1..2,
            geometry=-|
        ]
            {Съёмка}{Точка времени}

        \umlaggreg[
            mult1=1, 
            mult2=2,
            geometry=-|
        ]
            {Кадр}{Точка времени}


    \end{tikzpicture}
\end{document}

